\section{Por qué la estructura funcionó en su momento: traducción e Instruction Tuning como experimento natural}

La sección anterior completó la transformación formal; esta sección pregunta por la motivación histórica. La estructura adicional del encoder-decoder no era una complicación arbitraria: en 2017 la tarea central era la traducción automática, en la que la entrada y el objetivo provienen de idiomas distintos; más tarde, los datos académicos de instruction-tuning presentaban con frecuencia entradas largas y respuestas cortas. Cuando la estructura coincide con la distribución de los datos, el sesgo inductivo eleva de manera notable la eficiencia con un presupuesto limitado.

\subsection{Tres estructuras adicionales}

Esta subsección presenta primero la lista completa y después se centra en el primer elemento. En comparación con decoder-only, el encoder-decoder supone que la entrada y el objetivo son lo bastante distintos, que el objetivo debe leer una entrada ya codificada por completo y que los token de la entrada se prestan a una interacción totalmente bidireccional. Estos tres supuestos corresponden, respectivamente, a los parámetros, a las conexiones entre capas y a la attention mask.

\lecturefigure{hyung-slide-052.jpg}{Las tres estructuras adicionales del encoder-decoder respecto a decoder-only}{Hyung Won Chung official deck, p.~52}

\lecturefigure{hyung-slide-053.jpg}{Primer supuesto: la entrada y el objetivo difieren tanto que justifican parámetros independientes}{Hyung Won Chung official deck, p.~53}

\begin{importantbox}{Tres preguntas para juzgar si un inductive bias sigue siendo válido}
\begin{enumerate}
  \item ¿Con qué distribución de datos o restricción de la tarea coincidía originalmente?
  \item ¿Los datos, el objetivo y la forma de interacción actuales siguen cumpliendo esa restricción?
  \item Si se retira la estructura y se aumenta la escala, ¿en qué punto del presupuesto se cruzan las curvas de desempeño?
\end{enumerate}
Responder solo la primera pregunta lleva a confundir la validez histórica con una corrección permanente.
\end{importantbox}

\subsection{Traducción automática: los parámetros independientes eran naturales}

En 2017, el Transformer se orientaba a la source-to-target translation. La entrada es una oración en el idioma fuente y el objetivo es la salida en otro idioma; ambos difieren en vocabulario, orden de palabras y papel en la generación. En un escenario que solo optimiza la traducción, hacer que el encoder se especialice en el idioma fuente y el decoder en el idioma objetivo es un supuesto previo eficiente. Esta subsección reconoce primero esa validez histórica y después compara qué condiciones cambió el preentrenamiento general, para no deducir desde las interfaces de tareas actuales que el diseño anterior era «erróneo desde el principio».

\lecturefigure{hyung-slide-054.jpg}{En la traducción automática, la diferencia de distribución entre source y target respalda los parámetros independientes}{Hyung Won Chung official deck, p.~54}

\lecturefigure{hyung-slide-055.jpg}{Un modelo de lenguaje general debe integrar conocimiento entre idiomas, lo que debilita la motivación de los parámetros independientes}{Hyung Won Chung official deck, p.~55}

Al leer la segunda diapositiva conviene notar que la tarea ya cambió. El preentrenamiento moderno no solo aprende traducción del inglés al alemán, sino que debe formar un conocimiento del mundo compartido a partir de texto en muchos idiomas; un mismo hecho puede expresarse en idiomas distintos. Si la entrada y el objetivo quedan aislados de forma permanente, al modelo puede costarle más reutilizar sus representaciones. Por otro lado, la especialización multilingüe todavía puede tener valor, de modo que esta diapositiva plantea un scale-era counterargument, no que los parámetros independientes sean del todo inútiles.

\teachervoice{La comparación histórica de 00:27:27--00:28:45 es muy concreta: cuando el objetivo es solo traducir, separar los parámetros por idioma es «muy natural»; cuando el objetivo pasa a ser el aprendizaje de conocimiento general, el ponente prefiere fusionar el conocimiento entre idiomas. La evaluación de un inductive bias cambia junto con el objective.}

\subsection{FLAN: entradas largas y objetivos cortos dan una ventaja inesperada al encoder-decoder}

Instruction tuning ofrece un experimento natural más fino. El mismo equipo aplicó fine-tuning a PaLM (decoder-only) y a T5 (encoder-decoder); aunque dedicó la mayor parte del tiempo a optimizar PaLM, la ganancia de T5 fue mayor. El ponente no atribuyó el resultado a una superioridad general de la arquitectura, sino que volvió a la distribución de longitudes de los datos para buscar la correspondencia.

\lecturefigure{hyung-slide-056.jpg}{Instruction finetuning unifica una gran cantidad de tareas académicas como instrucciones en lenguaje natural}{Hyung Won Chung official deck, p.~56}

\lecturefigure{hyung-slide-057.jpg}{En el experimento FLAN, la ganancia del fine-tuning es mayor para el encoder-decoder}{Hyung Won Chung official deck, p.~57}

\begin{knowledgebox}{Lectura de la figura: se compara la ganancia, no qué modelo gana en términos absolutos}
La tabla destaca la improvement antes y después del instruction tuning. No basta con mirar el recuadro rojo para concluir que T5 supera a PaLM en todas las tareas y en todas las escalas; la conclusión de la clase es que la estructura del encoder-decoder coincide especialmente con ese conjunto de datos académicos, por lo que aprovecha mejor la señal del fine-tuning.
\end{knowledgebox}

\lecturefigure{hyung-slide-058.jpg}{Los datos académicos suelen tener una distribución de longitudes con entradas largas y objetivos cortos}{Hyung Won Chung official deck, p.~58}

Si la longitud de la entrada es la variable aleatoria $L_x$ y la longitud del objetivo es $L_y$, este conjunto de datos cumple aproximadamente
\[
\mathbb{E}[L_x]\gg \mathbb{E}[L_y].
\]
Aquí, $\mathbb{E}[L_x]$ y $\mathbb{E}[L_y]$ son las longitudes promedio de la entrada y del objetivo, respectivamente. Los enunciados largos y las etiquetas cortas hacen que las distribuciones de ambos lados difieran claramente, y los parámetros independientes de encoder/decoder pueden especializarse; pero la generación de textos largos y el chat de varios turnos reducen e incluso invierten esa asimetría.

\teachervoice{En 00:29:22--00:30:46, Hyung Won dice que el equipo dedicó el 99\% del tiempo a optimizar PaLM y, sin embargo, T5 obtuvo una ganancia mayor en pocos días; su explicación es que el sesgo de longitud de los datos y la arquitectura «coincidieron por casualidad». Este teacher voice convierte el resultado de una comparación de marcas en un análisis de correspondencia de distribuciones.}

\begin{warningbox}{La facilidad de evaluar un academic benchmark moldea la distribución de entrenamiento}
Los objetivos largos son difíciles de anotar manualmente y de calificar en forma automática, por lo que los datos académicos se inclinan hacia entradas largas y respuestas cortas. Que algo sea fácil de calificar no equivale a que tenga valor para el usuario: la escritura creativa, el código, el análisis y la conversación suelen requerir salidas largas. Si la arquitectura del modelo se optimiza solo para la distribución del benchmark, la ventaja estructural original puede desaparecer cuando cambian las tareas de despliegue.
\end{warningbox}

\subsection{Resumen de la sección}

La traducción automática y FLAN muestran por qué el sesgo inductivo puede producir beneficios reales en una época determinada: dirige una capacidad limitada hacia lo que la distribución de los datos más necesita. Pero el preentrenamiento general, la generación larga y el diálogo de varios turnos cambiaron la relación entre entrada y objetivo, y obligaron a los investigadores a volver a probar la separación de parámetros. La historia no sirve para ridiculizar los diseños anteriores, sino para identificar las condiciones en las que un diseño es válido.
